%!TEX root = ../template.tex
%%%%%%%%%%%%%%%%%%%%%%%%%%%%%%%%%%%%%%%%%%%%%%%%%%%%%%%%%%%%%%%%%%%%
%% chapter2.tex
%% NOVA thesis document file
%%
%% Chapter with the template manual
%%%%%%%%%%%%%%%%%%%%%%%%%%%%%%%%%%%%%%%%%%%%%%%%%%%%%%%%%%%%%%%%%%%%

\typeout{NT FILE chapter2.tex}%

\chapter{Theory of Computation}
\label{cha:theory_of_computation}


In the Theory of Computation, various models of computation with different levels of expressive power are studied. To discuss and compare different levels of expressive power, it is inevitable to consider algorithms in the context of some application domain. The typically chosen application domain is the syntax of formal languages.
There are two techniques for defining languages that are typically used. The first technique is based on the concept of generation and uses the formalism of grammars. Another technique is based on the mechanism of recognition and uses the formalism of automata.


\section{Chomsky Hierarchy}
\label{sec:chomsky_hierarchy}

Noam Chomsky (with the help of Marcel-Paul Schützenberger) identified several classes of formal languages, with different levels of specification complexity, requiring computational models with different levels of expressive power to solve the problem of recognizing these languages.

There are four levels: type 0, type 1, type 2, and type 3.

\begin{itemize}
  \item Type 0 — Recursively enumerable languages
  \item Type 1 — Context-sensitive languages
  \item Type 2 — Context-free languages
  \item Type 3 — Regular languages
  
\end{itemize}

The Chomsky hierarchy involves four levels of language types. Each type of language is associated with a type of generator and a type of recognizer.

\section{Context-free Languages}
\label{sec:context_free_languages}

Context-free languages include regular languages and are a type of formal language defined by context-free grammars. These grammars have production rules that generate sequences of symbols recursively, using non-terminal and terminal symbols. They have various applications in computer science, such as aiding in the implementation of compilers and describing document formats

\subsection{With a Local \LaTeX\ Installation} % (fold)
\label{sub:with_a_local_latex_installation}

